\scanpage{35}
\[
 \underline{\operatorname{Hom}}_{\mathrm{add}}(D,C')\xrightarrow{\simeq}
 \underline{\operatorname{Hom}}_{1/2\,\mathrm{add}}(\mathcal A,C')
\]
\unsure{d'où|allant} en sens inverse
\[
 \underline{\operatorname{Hom}}_{1/2\,\mathrm{add}}(\mathcal A,C')\xrightarrow{\simeq}
 \underline{\operatorname{Hom}}_{\mathrm{add}}(D,C')
\]
d'où un accouplement
\diagram{\begin{tikzcd}[ampersand replacement=\&,row sep=0pt,column sep=0.6em]
D \& \times \& \underline{\operatorname{Hom}}_{1/2\,\mathrm{add}}(\mathcal A,C') \arrow[r] \& C' \\
\rotatebox{-90}{$\subset$} \& \& \& \\
\underline{\operatorname{Hom}}_{1/2\,\mathrm{add}}(\mathcal A^\circ,\mathrm{Ab}) \& \& \&
\end{tikzcd}}
qu'on note par \deleted{\unsure{la|ex.}}
\[
 (F,G)\unsure{\mapsto|\vdash} F\underset{C}{\otimes}G
\]
[dans le cas où $\mathcal A$ est réduit à un
élément, donc où sa str. $1/2$ add. est
\unsure{défini|définie} par un anneau $k$,
$F$ est un \inserted{$k$-}module à droite libre de
t. f, $G$ est un \deleted{\unsure{$k$}}objet de $C'$ où $k$
opère à g (i.e. un $k$-$C'$-module
à gauche) et $F\otimes G$ est le
produit tensoriel ``ordinaire''.

Si on considère \deleted{\unsure{ce|alors}}le plongement
dual $C\xrightarrow{\sim}D'{}^\circ\subset\underline{\operatorname{Hom}}_{1/2\,\mathrm{add}}(\mathcal A,\mathrm{Ab})^\circ$,
on trouve
\diagram{\begin{tikzcd}[ampersand replacement=\&]
\underline{\operatorname{Hom}}_{\mathrm{add}}(D'{}^\circ,C') \arrow[r,"{\simeq}"] \&
\underline{\operatorname{Hom}}_{1/2\,\mathrm{add}}(\mathcal A,C') \arrow[l,bend left=40,"{\simeq}"]
\end{tikzcd}}
d'où
% archive sheet number: 34
